% !TEX root = ../../../main.tex
% Week 4, IMG_9238 and IMG_9240--IMG_9242. Keep the syllabus location.
% Only a final sentence is visible in IMG_9242; a hand obscures the rest.
\section{Power series}
\label{sec:analysis-i:power-series}

A power series uses a variable as the base of each term. For each fixed
value of that variable, it is an ordinary number series, so the tests
from \autoref{sec:analysis-i:comparison-tests} apply.

\begin{definition}
  A \emph{power series centred at zero} has the form
  \[
    \sum_{n=0}^{\infty}c_nx^n,
  \]
  where the coefficients $c_n$ are fixed real numbers.
  Its \emph{radius of convergence} is a number $R\in[0,+\infty]$
  such that the series converges absolutely for $|x|<R$ and diverges
  for $|x|>R$. The points $x=\pm R$, when $0<R<\infty$, must be
  examined separately.
\end{definition}

% Addition awaiting author review. The source introduces
% radii through examples; this supplies the existence assertion used above.
\begin{pendingreview}
\begin{proposition}
  \label{thm:analysis-i:power-series-radius}
  Every real power series centred at zero has a radius of convergence
  $R\in[0,+\infty]$ with the properties in the preceding definition.
\end{proposition}

\begin{proof}
  At $x=0$ the constant term is $c_0$ and every other term vanishes,
  so the series converges there. Suppose it converges at some
  $x_0\ne0$. Its terms tend to zero by
  \autoref{thm:analysis-i:series-term-test}, so there is a number $B>0$
  such that $|c_nx_0^n|\le B$ for every $n\ge0$. If $|x|<|x_0|$, then
  \[
    |c_nx^n|\le B\left(\frac{|x|}{|x_0|}\right)^n.
  \]
  The geometric majorant converges. Thus convergence at any nonzero
  point implies absolute convergence at every point closer to zero.

  Let $S$ consist of the nonnegative values $r$ where the power series
  converges absolutely. This set contains $0$.
  If $S$ is bounded above, let $R=\sup S$; otherwise let $R=+\infty$.
  If $|x|<R$, there is an $r\in S$ with $|x|<r$. The preceding
  comparison then proves absolute convergence at $x$. For $R=0$
  this interior condition is empty, while convergence at zero remains.

  If $R<\infty$ and the series converged at a point with $|x|>R$,
  choose $t$ with $R<t<|x|$. The same comparison would give absolute
  convergence at $t$, so $t\in S$, contradicting that $R$ is an upper
  bound. The series therefore diverges for $|x|>R$. This proves all
  the required properties. Boundary behaviour when $0<R<\infty$
  still requires separate tests.
\end{proof}
\end{pendingreview}

An infinite radius means absolute convergence for every real $x$.

\subsection{The geometric power series}

The series $\sum_{n=0}^{\infty}x^n$ has radius $R=1$, since
\autoref{sec:analysis-i:number-series} proves that it converges exactly
for $|x|<1$. Neither endpoint is included. At $x=1$ every term is $1$;
at $x=-1$ the terms alternate between $1$ and $-1$. In both cases the
terms fail to tend to zero.
% IMG_9238's annotation R=|x| confuses the term ratio with the radius.

\subsection{The exponential series}

Consider the polynomials and their limiting series
\[
  E_N(x)=\sum_{k=0}^{N}\frac{x^k}{k!},\qquad
  E(x)=\sum_{k=0}^{\infty}\frac{x^k}{k!}.
\]
For fixed $x\ne0$, the absolute successive-term ratio is
\[
  \frac{|x|^{k+1}/(k+1)!}{|x|^k/k!}
  =\frac{|x|}{k+1}\longrightarrow0.
\]
The ratio test proves absolute convergence for every such $x$.
At $x=0$, only the constant term remains, so $E(0)=1$.
Thus this series has infinite radius of convergence and $E_N(x)\to E(x)$
for every fixed real $x$. The lecture identifies $E(x)$ with $e^x$;
the properties of the exponential function are not proved by this
convergence argument. At $x=1$ the sum is the earlier factorial-series
definition of $e$.

\subsection{A family with different endpoints}

Let $s\ge0$ be a real number and consider
\[
  F_s(x)=\sum_{n=1}^{\infty}\frac{x^n}{n^s}.
\]
As in the $p$-series example, real powers use their usual exponent and
order laws. For $x\ne0$, the $n$th root of the absolute value of the
$n$th term is
\[
  \left|\frac{x^n}{n^s}\right|^{1/n}
  =\frac{|x|}{(n^{1/n})^s}\longrightarrow |x|.
\]
To justify the denominator's limit using the sequence results already
available, choose an integer $m\ge s$. Since $n^{1/n}\ge1$,
\[
  1\le(n^{1/n})^s\le(n^{1/n})^m\longrightarrow1.
\]
The last limit follows from \autoref{thm:analysis-i:root-sequence-limits}
and the product rule, and the squeeze principle gives the desired limit.

For $|x|<1$, the root test proves absolute convergence, with $x=0$
immediate. For $|x|>1$, choose $1<r<|x|$. The root values eventually
exceed $r$, so $|x^n/n^s|>r^n$ for every sufficiently large $n$.
The terms do not tend to zero and the series diverges. Therefore $R=1$
for every $s\ge0$.

\begin{marginfigure}
  \centering
  % !TEX root = ../../main.tex
% Shared physical dimensions for interval and number-line diagrams.
\tikzset{
  interval diagram/.style={
    x=1cm,y=6mm,font=\footnotesize,
    color=black,line width=0.5pt,line cap=round,line join=round,
    >={Stealth[length=4pt,width=3pt]}
  },
  interval endpoint/.style={
    circle,draw=black,line width=0.5pt,
    minimum size=4pt,inner sep=0pt,outer sep=0pt
  },
  interval open/.style={interval endpoint,fill=white},
  interval closed/.style={interval endpoint,fill=black},
  interval label/.style={inner sep=1pt}
}

\begin{tikzpicture}[interval diagram,x=13mm,y=13mm]
  \foreach \row/\case in {2/{s=0},1/{0<s\le1},0/{s>1}}{
    \draw[gray,<->] (-1.6,\row) -- (1.6,\row);
    \draw (-1,\row) -- (1,\row);
    \node[above=9pt] at (0,\row) {$\case$};
    \node[below=6pt] at (-1,\row) {$-1$};
    \node[below=6pt] at (1,\row) {$1$};
  }
  \node[interval open] at (-1,2) {};
  \node[interval open] at (1,2) {};
  \node[interval closed] at (-1,1) {};
  \node[interval open] at (1,1) {};
  \node[interval closed] at (-1,0) {};
  \node[interval closed] at (1,0) {};
\end{tikzpicture}

  \caption[Three endpoint patterns for radius one.]{All three cases have radius $1$. Open or filled endpoint
    markers record the separate tests at $-1$ and $1$; the radius
    alone does not settle them.}
  \label{fig:power-series-endpoints}
\end{marginfigure}

At $x=1$, this is the $p$-series, which converges exactly when $s>1$.
At $x=-1$, it is $\sum(-1)^n/n^s$, which converges for $s>0$ by
\autoref{ex:analysis-i:alternating-power-series}. For $s=0$, the terms
do not tend to zero. Accordingly the intervals of convergence are
\[
  \begin{array}{c|c}
    s=0 & (-1,1)\\[3pt]
    0<s\le1 & [-1,1)\\[3pt]
    s>1 & [-1,1].
  \end{array}
\]
\autoref{fig:power-series-endpoints} keeps the three cases distinct.
% IMG_9242 claims [-1,1] without a visible condition. The right endpoint
% requires s>1; the left requires s>0. No obscured writing is transcribed.

The following exercise introduces a scale factor in the coefficients.
Use the series tests for the interior and exterior, then examine
absolute convergence at each endpoint separately.

% !TEX root = ../../hw04.tex
% Source: HW4 Intro to Math Analysis I.pdf, page 2, question 7.
% The integer k is unrestricted, so zero and negative integers are included.
\begin{exercise}
  \label{ex:hw04:scaled-power-series}
  \solutionlink{hw04:scaled-power-series}
  For $x\in\R$, investigate the absolute convergence of the power series
  \[
    \sum_{n=1}^{\infty}\frac{x^n}{(\sqrt3)^n n^k},
  \]
  where $k$ is an integer using the ratio test or root test for convergence
  or divergence.
\end{exercise}

